\documentclass[11pt,letterpaper]{article}
\usepackage[utf8]{inputenc}
\usepackage[T1]{fontenc}
\usepackage[spanish,es-noshorthands]{babel}
\usepackage{lmodern}
\usepackage[margin=2.3cm]{geometry}
\usepackage{amsmath,amssymb,amsthm}
\usepackage{booktabs,tabularx,array}
\usepackage{xcolor}
\usepackage{listings}
\usepackage{enumitem}
\usepackage{fancyhdr}
\usepackage[hidelinks]{hyperref}
\setlength{\headheight}{14pt}
\sloppy

\definecolor{fondo}{RGB}{246,246,246}
\definecolor{coment}{RGB}{90,120,90}
\lstset{
  basicstyle=\ttfamily\footnotesize,
  backgroundcolor=\color{fondo},
  commentstyle=\color{coment}\itshape,
  keywordstyle=\bfseries,
  frame=single, framerule=0.3pt, rulecolor=\color{gray},
  breaklines=true, columns=fullflexible, keepspaces=true,
  showstringspaces=false, upquote=true,
  literate={ñ}{{\~n}}1 {á}{{\'a}}1 {é}{{\'e}}1 {í}{{\'i}}1 {ó}{{\'o}}1 {ú}{{\'u}}1
           {Þ}{{\TH}}1 {É}{{\'E}}1 {³}{{$^3$}}1 {Ó}{{\'O}}1 {ê}{{\^e}}1
}
\lstdefinestyle{salida}{language={}, numbers=none}
\lstdefinestyle{py}{language=Python, numbers=left, numberstyle=\tiny\color{gray}}

\newcommand{\xor}{\oplus}
\newcommand{\B}{\{0,1\}^{128}}
\newcommand{\paso}[1]{\par\smallskip\noindent\textbf{#1.}\ }
\newcommand{\tesis}[1]{\par\smallskip\noindent\fbox{\parbox{0.96\linewidth}{\textbf{Tesis.} #1}}\par}

\newenvironment{declaracion}{\par\medskip\noindent\textbf{Declaración de variables.}\begin{itemize}[leftmargin=1.5em,itemsep=1pt]}{\end{itemize}}

\pagestyle{fancy}
\fancyhf{}
\lhead{Criptografía simétrica}
\rhead{Taller B --- Grupo Blockchain}
\cfoot{\thepage}

\title{\textbf{Taller B: El registro}\\[2pt]
\large Qué hace que una cadena de bloques no se pueda reescribir}
\author{Integrantes: \underline{\hspace{3.5cm}}\quad \underline{\hspace{3.5cm}}\\[4pt]
\underline{\hspace{3.5cm}}\quad \underline{\hspace{3.5cm}}\\[6pt]
Profesor: Andrey Lizarazo}
\date{Fecha: \underline{\hspace{3cm}}}

\begin{document}
\maketitle
\thispagestyle{fancy}

\section*{Resumen del caso}

Un usuario acusa al mercado de licencias de haberle borrado una compra; el mercado responde que
eso es imposible porque ``la cadena verifica''. En este informe construimos el mismo registro de
dos maneras ---con CBC y con Davies--Meyer--- y mostramos que:
\begin{itemize}[itemsep=1pt]
  \item la cadena CBC es una biyección conocida la clave: quien la tiene reescribe cualquier
        historia en microsegundos y el verificador la acepta; peor aún, borrar la primera o la
        última transferencia ni siquiera requiere la clave;
  \item la cadena Davies--Meyer no se deja retroceder: el ataque se detiene en una ecuación
        $E_{m}(x)\xor x = h$ que no admite despeje, y que solo sirve si el sello de referencia
        está publicado fuera del control del mercado.
\end{itemize}
Conclusión del caso: \emph{el argumento del mercado es falso}. Que la cadena CBC verifique no prueba
nada sobre si el mercado la alteró. El usuario tiene razón en que el borrado era posible y además indetectable
para el verificador.

\section*{Declaración general de variables}

Estas declaraciones valen para todo el informe; antes de cada ejercicio se repiten las que
intervienen en él.
\begin{itemize}[itemsep=1pt]
  \item $n = 128$: número de bits de un bloque. $\B$: conjunto de bloques de 128 bits (16 bytes).
  \item $\mathcal{K} = \B$: espacio de claves de AES-128.
  \item $E:\mathcal{K}\times\B\to\B$, $(k,x)\mapsto E_k(x)$: cifrado AES. Para cada $k$ fija,
        $E_k$ es una permutación de $\B$ y $D_k = E_k^{-1}$, es decir
        $D_k(E_k(x)) = x$ y $E_k(D_k(y)) = y$ para todo $x,y\in\B$.
  \item $\xor:\B\times\B\to\B$: suma XOR bit a bit. Propiedades que se citarán por nombre:
    \begin{itemize}[itemsep=0pt]
      \item[(X1)] \emph{autoinversa}: $x\xor x = 0$;
      \item[(X2)] \emph{neutro}: $x \xor 0 = x$;
      \item[(X3)] \emph{asociativa y conmutativa}.
    \end{itemize}
  \item $t\in\mathbb{Z}$, $t\ge 1$: número de entradas del registro (aquí $t=4$).
  \item $i\in\{1,\dots,t\}$: índice de bloque. $m_i\in\B$: entrada $i$ en claro.
        $c_i\in\B$: bloque cifrado $i$. $h_i\in\B$: valor acumulado de la cadena tras el bloque $i$.
  \item $\|$: concatenación de bloques.
\end{itemize}

Todo el código está en la carpeta \texttt{codigo/}: \texttt{registro\_tool.py} (las tres funciones
pedidas) y \texttt{experimentos.py} (las mediciones y ataques). Las salidas reproducidas en este
informe son literales.

%======================================================================
\section{Ejercicio 1. La recurrencia y el error \hfill (25 puntos)}

\begin{declaracion}
  \item $k\in\mathcal{K}$: clave del mercado, \texttt{4ac43a52cc764438cae7b3af9fced5a9}.
  \item $IV\in\B$: valor de arranque, \texttt{3c1cc2d73a499ca72f36e7303d76a130}. Convenio $c_0 := IV$.
  \item $m_1,\dots,m_4\in\B$ y $c_1,\dots,c_4\in\B$, ligados por
        $c_i = E_k(m_i\xor c_{i-1})$.
  \item $e\in\B$, $e\neq 0$: máscara de error introducida en un bloque cifrado.
  \item $c_3' = c_3\xor e\in\B$: bloque $c_3$ alterado; $m_3', m_4'\in\B$: entradas descifradas con el bloque alterado.
\end{declaracion}

\subsection*{(a) Leer el registro}

\paso{Despeje de la dirección inversa}
Partimos de la ecuación del enunciado, $c_i = E_k(m_i\xor c_{i-1})$. Aplicamos $D_k$ a ambos
lados y usamos $D_k\circ E_k=\mathrm{id}$:
\[
  D_k(c_i) = m_i \xor c_{i-1}.
\]
Sumamos $c_{i-1}$ a ambos lados; por (X3) y (X1), $c_{i-1}\xor c_{i-1}=0$, y por (X2):
\[
  D_k(c_i)\xor c_{i-1} = m_i\xor (c_{i-1}\xor c_{i-1}) = m_i \xor 0 = m_i .
\]
Así, la recurrencia en sentido inverso es $m_i = D_k(c_i)\xor c_{i-1}$, con $c_0=IV$.

\paso{Implementación} Las dos direcciones, tal como quedaron en \texttt{registro\_tool.py}:
\begin{lstlisting}[style=py]
def cbc_descifrar(k, iv, cadena):
    entradas, anterior = [], iv
    for c in bloques(cadena):
        entradas.append(xor(aes_d(k, c), anterior))   # m_i = D_k(c_i) XOR c_{i-1}
        anterior = c
    return entradas

def cbc_cifrar(k, iv, entradas):
    salida, anterior = b"", iv
    for m in entradas:
        assert len(m) == 16
        anterior = aes_e(k, xor(m, anterior))         # c_i = E_k(m_i XOR c_{i-1})
        salida += anterior
    return salida
\end{lstlisting}

\paso{Resultado} \texttt{python3 registro\_tool.py leer} devuelve:

\begin{center}
\begin{tabular}{clll}
\toprule
$i$ & $c_i$ & $m_i$ (hex) & $m_i$ (texto) \\
\midrule
1 & \texttt{acc42eca\ldots a91947} & \texttt{545830313a4c303034323a45512d4120} & \texttt{TX01:L0042:EQ-A\textvisiblespace} \\
2 & \texttt{bcba08d8\ldots b4781f} & \texttt{545830323a4c303034323a45512d4220} & \texttt{TX02:L0042:EQ-B\textvisiblespace} \\
3 & \texttt{7d03a386\ldots 74c0cd} & \texttt{545830333a4c303137373a45512d4120} & \texttt{TX03:L0177:EQ-A\textvisiblespace} \\
4 & \texttt{8210d2b6\ldots a52261} & \texttt{545830343a4c303137373a45512d4220} & \texttt{TX04:L0177:EQ-B\textvisiblespace} \\
\bottomrule
\end{tabular}
\end{center}

Es decir: la licencia L0042 pasó a EQ-A (TX01) y luego a EQ-B (TX02); la L0177 pasó a EQ-A (TX03)
y luego a EQ-B (TX04). Además comprobamos la ida y vuelta:
\\ \texttt{cbc\_cifrar(k, iv, cbc\_descifrar(k, iv, c)) == c} da \texttt{True}.
El verificador del mercado (\texttt{python3 verificador.py cbc registro\_cbc.json}) imprime
exactamente las mismas cuatro entradas y responde \texttt{CADENA VALIDA}.

\subsection*{(b) Romper un byte}

\paso{1. Experimento} Alteramos el primer byte de $c_3$ con la máscara
$e = \texttt{40}\,\texttt{00}\cdots\texttt{00}$ (un solo bit) y desciframos de nuevo
(\texttt{python3 registro\_tool.py romper 3} y \texttt{experimentos.py}). Comparación byte a byte:

\begin{center}
\small
\begin{tabular}{clll}
\toprule
Entrada & Original & Tras el daño & Diferencia ($m_i\xor m_i'$) \\
\midrule
1 & \texttt{545830313a4c\ldots2d4120} & \texttt{545830313a4c\ldots2d4120} & 0 bytes, 0 bits: \textbf{intacta}\\
2 & \texttt{545830323a4c\ldots2d4220} & \texttt{545830323a4c\ldots2d4220} & 0 bytes, 0 bits: \textbf{intacta}\\
3 & \texttt{545830333a4c\ldots2d4120} & \texttt{de2854c98577\ldots57eaad} & 16 bytes, 68 bits: \textbf{destruida}\\
4 & \texttt{545830343a4c\ldots2d4220} & \texttt{145830343a4c\ldots2d4220} & 1 byte, 1 bit (\texttt{0x40}): \textbf{alterada}\\
\bottomrule
\end{tabular}
\end{center}

La diferencia completa de la entrada 3 es
\texttt{8a7064fabf3ba44a683e6bf6827aab8d}: no guarda ninguna relación visible con $e$.
La de la entrada 4 es \texttt{40000000000000000000000000000000}, exactamente $e$:
la \texttt{T} (\texttt{0x54}) pasó a \texttt{0x14} y nada más cambió.
Con este registro dañado el verificador responde \texttt{CADENA INVALIDA} y marca las entradas 3 y 4.

\paso{Por qué (derivación)} Con $c_3' = c_3\xor e$ y los demás bloques sin tocar:
\begin{itemize}[itemsep=2pt]
  \item $m_1, m_2$: sus fórmulas $m_1 = D_k(c_1)\xor IV$ y $m_2 = D_k(c_2)\xor c_1$ no contienen a $c_3$; quedan intactas.
  \item $m_3' = D_k(c_3\xor e)\xor c_2$. El error entra \textbf{por dentro} de $D_k$. Como $D_k$ es una permutación
        pseudoaleatoria, $D_k(c_3\xor e)$ no tiene relación útil con $D_k(c_3)$: cada bit de salida cambia con
        probabilidad $\approx 1/2$ (observamos 68 de 128). La entrada queda destruida e impredecible para quien no tenga $k$.
  \item $m_4' = D_k(c_4)\xor c_3' = D_k(c_4)\xor c_3\xor e = m_4\xor e$ por (X3). El error entra \textbf{por fuera},
        como sumando del XOR final, y pasa intacto: cambian exactamente los bits de $e$.
  \item $m_5,\dots$ (si existieran) solo dependen de $c_4, c_5,\dots$ y quedarían intactas: el daño se propaga a
        dos bloques y se detiene.
\end{itemize}

\paso{2. La asimetría y su consecuencia práctica}
En $m_3$ el error está \emph{dentro} de $D_k$; en $m_4$ está \emph{fuera}. Para quien daña a propósito
la diferencia es total:
\begin{itemize}[itemsep=2pt]
  \item Sobre $m_4$ tiene \textbf{control exacto y sin clave}: si quiere que $m_4$ pase a valer $m_4^*$, le basta
        usar $e = m_4\xor m_4^*$ (es la ecuación $a\xor x=b$ del Taller 0). Lo comprobamos: con
        $e = (\texttt{B}\xor\texttt{A})$ en la posición 14 de $c_3$, la entrada 4 pasa a decir
        \texttt{TX04:L0177:EQ-A}, sin conocer $k$.
  \item Sobre $m_3$ no tiene ningún control: el precio de su manipulación es que la entrada anterior
        sale como basura. En este registro esa basura es detectable porque no cumple el formato
        (la probabilidad de que 16 bytes aleatorios cumplan
        \verb|TX\d\d:L\d\d\d\d:[A-Z][A-Z]-[A-Z] | es $100\cdot 10^4\cdot 26^3/2^{128}\approx 5\cdot10^{-29}$).
\end{itemize}
Es decir: CBC es \emph{maleable}. El adversario sin clave puede dirigir bits de un bloque a costa de
destruir el anterior. Lo único que lo delata es la redundancia del formato, no el encadenamiento.
Y si el bloque que se sacrifica es el IV (que viaja en claro), ni siquiera hay basura: con
$IV' = IV\xor m_1\xor m_1^*$ se cambia $m_1$ a voluntad y todo lo demás queda intacto.

%======================================================================
\section{Ejercicio 2. Reescribir la historia \hfill (30 puntos)}

\begin{declaracion}
  \item $k\in\mathcal{K}$, $IV\in\B$ como en el ejercicio 1; $t\in\mathbb{Z}$, $t\ge 1$.
  \item $M = (\B)^t$: conjunto de historias $(m_1,\dots,m_t)$. $C=(\B)^t$: conjunto de cadenas $(c_1,\dots,c_t)$.
  \item $F_{k,IV}: M\to C$, la aplicación de cifrado CBC; $G_{k,IV}: C\to M$, la candidata a inversa.
  \item $m^* = (m_1^*,\dots,m_t^*)\in M$: historia alternativa (``lo que quiero que diga'').
\end{declaracion}

\subsection*{(a) Hacerlo}

Historia falsa: la misma, salvo que TX02 dice que L0042 se quedó en EQ-A:
\[
m^* = (\texttt{TX01:L0042:EQ-A\textvisiblespace},\ \texttt{TX02:L0042:EQ-A\textvisiblespace},\ \texttt{TX03:L0177:EQ-A\textvisiblespace},\ \texttt{TX04:L0177:EQ-B\textvisiblespace}).
\]
Con \texttt{python3 registro\_tool.py reescribir} se cifra completa con \texttt{cbc\_cifrar} y se
guarda en \texttt{falsificado.json}. Cadena resultante:

\begin{center}
\small
\begin{tabular}{cll}
\toprule
$i$ & $c_i^*$ (falsificada) & $c_i$ (original) \\
\midrule
1 & \texttt{acc42ecaebdb06d4e3b2852697a91947} & \texttt{acc42ecaebdb06d4e3b2852697a91947} \\
2 & \texttt{8f3ebfc015f257f20935a4a9bb070b22} & \texttt{bcba08d852933bcb8614195bedb4781f} \\
3 & \texttt{3314b26facf1eaae43ea1b86beb9ec3b} & \texttt{7d03a3865c756f1ed89c5c56fd74c0cd} \\
4 & \texttt{0b1fc3366cb2c953363114fc04286dc5} & \texttt{8210d2b695c9514852ac55dcefa52261} \\
\bottomrule
\end{tabular}
\end{center}

$c_1$ no cambia porque $m_1$ no cambió; desde $c_2$ todo cambia, porque cada bloque arrastra al
siguiente. Esto \emph{no} es una protección: quien tiene la clave simplemente recalcula todo.

\paso{Tiempo} Promedio de 1000 reescrituras completas: \textbf{22 microsegundos} por historia
(medido con \texttt{time.perf\_counter} en \texttt{experimentos.py}).

\paso{Respuesta literal del verificador} \texttt{python3 verificador.py cbc falsificado.json}:
\lstinputlisting[style=salida]{../codigo/salida_verificador_cbc_falsificado.txt}
El verificador no se dio cuenta de nada: acepta la historia falsa con el mismo mensaje con el que
acepta la verdadera.

\subsection*{(b) Demostración: la reescritura siempre funciona}

\paso{1}
\tesis{Fijados $k\in\mathcal{K}$ e $IV\in\B$, la aplicación $F_{k,IV}:(\B)^t\to(\B)^t$,
$(m_1,\dots,m_t)\mapsto(c_1,\dots,c_t)$ con $c_0=IV$ y $c_i=E_k(m_i\xor c_{i-1})$, es biyectiva, y su
inversa es $G_{k,IV}(c_1,\dots,c_t) = (m_1,\dots,m_t)$ con $m_i = D_k(c_i)\xor c_{i-1}$, $c_0=IV$.}

\paso{Hipótesis}
(H1) $k\in\mathcal{K}$ e $IV\in\B$ son fijos y conocidos.
(H2) Para $k$ fija, $E_k$ es una permutación de $\B$ con inversa $D_k$:
$D_k(E_k(x))=x$ y $E_k(D_k(y))=y$ para todo $x,y\in\B$.
(H3) $(\B,\xor)$ cumple (X1), (X2), (X3).
(H4) $t\ge1$, y $F_{k,IV}$, $G_{k,IV}$ están definidas por las fórmulas de la tesis.
Ambas están bien definidas: cada componente se calcula a partir de datos ya disponibles
($c_{i-1}$ es dato de entrada en $G$, y en $F$ se calcula en el paso anterior).

\paso{Planteamiento}
Una aplicación entre dos conjuntos es biyectiva si y solo si tiene inversa a ambos lados. Probaremos
(i) $G\circ F = \mathrm{id}_M$ y (ii) $F\circ G = \mathrm{id}_C$. Partimos de la ecuación del enunciado,
$c_i = E_k(m_i\xor c_{i-1})$.

\paso{Comparación}
\emph{(i) $G\circ F=\mathrm{id}_M$.} Sea $m=(m_1,\dots,m_t)\in M$ y $c = F(m)$. Sea $(\tilde m_1,\dots,\tilde m_t) = G(c)$.
Para cada $i$, con $c_0=IV$ en ambas aplicaciones:
\begin{align*}
\tilde m_i &= D_k(c_i)\xor c_{i-1} && \text{(definición de } G)\\
&= D_k\bigl(E_k(m_i\xor c_{i-1})\bigr)\xor c_{i-1} && \text{(definición de } F)\\
&= (m_i\xor c_{i-1})\xor c_{i-1} && \text{(H2: } D_k\circ E_k=\mathrm{id})\\
&= m_i\xor (c_{i-1}\xor c_{i-1}) && \text{(X3: asociatividad)}\\
&= m_i \xor 0 && \text{(X1: } x\xor x=0)\\
&= m_i && \text{(X2: } x\xor 0 = x).
\end{align*}
Observamos que $G$ usa los $c_{i-1}$ \emph{de la cadena} y $F$ usa los $c_{i-1}$ que ella misma produjo: son los mismos
porque $c=F(m)$. Por tanto $G(F(m))=m$.

\emph{(ii) $F\circ G=\mathrm{id}_C$.} Sea $c=(c_1,\dots,c_t)\in C$, $m = G(c)$ y $(\tilde c_1,\dots,\tilde c_t) = F(m)$, con
$\tilde c_0 = c_0 = IV$. Probamos por inducción sobre $i$ que $\tilde c_i = c_i$.
Caso base $i=0$: $\tilde c_0 = IV = c_0$. Paso inductivo: supongamos $\tilde c_{i-1}=c_{i-1}$. Entonces
\begin{align*}
\tilde c_i &= E_k(m_i\xor \tilde c_{i-1}) && \text{(definición de } F)\\
&= E_k(m_i\xor c_{i-1}) && \text{(hipótesis inductiva)}\\
&= E_k\bigl(D_k(c_i)\xor c_{i-1}\xor c_{i-1}\bigr) && \text{(definición de } G)\\
&= E_k\bigl(D_k(c_i)\xor 0\bigr) = E_k(D_k(c_i)) && \text{(X3, X1, X2)}\\
&= c_i && \text{(H2: } E_k\circ D_k=\mathrm{id}).
\end{align*}
Luego $F(G(c)) = c$.

\paso{Cierre lógico}
$F$ tiene inversa por la izquierda (i), luego es inyectiva; tiene inversa por la derecha (ii), luego es
sobreyectiva. Una aplicación inyectiva y sobreyectiva es biyectiva, y su inversa es única e igual a $G$.
La demostración no usa ninguna propiedad de los $m_i$: vale para \emph{toda} historia de longitud $t$.

\paso{Conclusión}
Conocidos $k$ e $IV$, $F_{k,IV}$ es una biyección de $(\B)^t$ en sí mismo con inversa explícita
$G_{k,IV}$, $m_i = D_k(c_i)\xor c_{i-1}$. \hfill$\blacksquare$

\paso{2. Consecuencia}
Sí: por la sobreyectividad (en realidad basta con que $F$ esté definida en todo $M$), para toda historia
alternativa $m^*\in M$ existe una cadena que la verifica, a saber $c^* = F_{k,IV}(m^*)$, y por la inyectividad es la única
con ese $IV$. Para construirla hace falta exactamente lo mismo que para cifrar: la clave $k$ y el $IV$
(que es público). No hace falta ninguna búsqueda; es un cálculo directo de $t$ cifrados AES, que
medimos en 22\,$\mu$s. Como el ``verificador'' solo comprueba que $G(c^*)$ tenga un formato válido, y
$G(c^*) = m^*$ tiene formato válido por construcción, acepta $c^*$ siempre.

\subsection*{(c) Dictamen}

\paso{Contra qué sí protege}
Contra la \emph{corrupción accidental} de los bloques cifrados: un error aleatorio en $c_j$ destruye
$m_j$ (ejercicio 1) y, con probabilidad $1-5\cdot10^{-29}$, la basura no cumple el formato y el
verificador lo detecta. También impide que un tercero \emph{sin} la clave lea el registro: es
confidencialidad, no integridad.

\paso{Contra qué no protege nunca}
\begin{itemize}[itemsep=2pt]
  \item Contra la \emph{alteración deliberada} por quien tiene la clave: es decir, el propio mercado, que es
        quien custodia el registro. Por (b), cualquier historia se reescribe en microsegundos y verifica.
  \item Contra un tercero sin clave que \emph{borra extremos}. Lo comprobamos
        (\texttt{sin\_ultima.json}, \texttt{sin\_primera.json}): quitar $c_4$ deja una cadena de tres
        entradas válidas; quitar $c_1$ y poner $IV' = c_1$ deja TX02, TX03, TX04 válidas. En ambos casos el
        verificador responde \texttt{CADENA VALIDA} sin que se haya usado la clave. Es precisamente
        ``borrar una compra''.
  \item Contra la manipulación dirigida de bits sin clave (maleabilidad, ejercicio 1(b)), salvo por la
        redundancia del formato.
\end{itemize}

\paso{Quién tenía razón}
El mercado no. Su argumento ---``la cadena verifica, luego no se pudo alterar''--- es falso: hemos
construido una cadena alterada que verifica, y hemos demostrado que eso es posible para cualquier
historia. Que la cadena verifique es compatible tanto con que el mercado sea honesto como con que haya
borrado la compra. Con este registro el usuario no puede \emph{probar} que se la borraron, pero el
mercado tampoco puede probar que no lo hizo; y la afirmación de que ``es imposible'' es falsa.

%======================================================================
\section{Ejercicio 3. La cadena que no se deja \hfill (45 puntos)}

\begin{declaracion}
  \item $h_0 = 0^{128}\in\B$: valor inicial público. $h_1,\dots,h_4\in\B$: valores acumulados.
  \item $m_1,\dots,m_4\in\B$: entradas en claro, \emph{públicas}, que en Davies--Meyer hacen de clave:
        $m_i\in\mathcal{K}$.
  \item $f:\mathcal{K}\times\B\to\B$, $f(m,h) = E_m(h)\xor h$: función de compresión de Davies--Meyer.
        Recurrencia: $h_i = f(m_i,h_{i-1}) = E_{m_i}(h_{i-1})\xor h_{i-1}$.
  \item $g:\mathcal{K}\times\B\to\B$, $g(m,h)=E_m(h)$: la misma construcción sin el término final.
  \item $S\in\B$: sello publicado, \texttt{51cfb2a7bf6b973b899c232b4ecb521c}.
  \item $m_2^* = \texttt{TX02:L0042:EQ-A\textvisiblespace}$; $h_i^*$: valores acumulados con la historia alterada.
  \item $x\in\B$: incógnita.
\end{declaracion}

\subsection*{(a) Construir y explorar}

\paso{1. Implementación}
\begin{lstlisting}[style=py]
def davies_meyer(h0, entradas):
    h = h0
    for m in entradas:
        assert len(m) == 16
        h = xor(aes_e(m, h), h)                       # el mensaje va como clave
    return h
\end{lstlisting}
El orden es \texttt{aes\_e(m, h)}: el bloque de mensaje es la clave de AES y el acumulado es el dato.
Valores intermedios, idénticos a los que imprime \texttt{python3 verificador.py dm registro\_dm.json}:

\begin{center}
\begin{tabular}{cl}
\toprule
$i$ & $h_i = E_{m_i}(h_{i-1})\xor h_{i-1}$\\
\midrule
0 & \texttt{00000000000000000000000000000000}\\
1 & \texttt{692d1a9c017dc9fea6251bf243385b4b}\\
2 & \texttt{e0e9a82bb76575a94a4f0153be97c127}\\
3 & \texttt{460eaf9791f61ab80f21f1601665fcb7}\\
4 & \texttt{51cfb2a7bf6b973b899c232b4ecb521c}\\
\bottomrule
\end{tabular}
\end{center}

$h_4$ coincide con el sello de \texttt{sello\_publicado.txt}: \texttt{python3 registro\_tool.py sello}
responde \texttt{COINCIDE}. El sello vive fuera de \texttt{registro\_dm.json} porque es la
referencia contra la que se compara el registro: si estuviera en el mismo archivo, quien edita el archivo
editaría también la referencia (ver (c)).

\paso{2. El mismo ataque del ejercicio 2}
Cambiamos $m_2$ por $m_2^*$ y recomputamos (\texttt{dm\_falsificado.json}):

\begin{center}
\begin{tabular}{cll}
\toprule
$i$ & $h_i^*$ (historia alterada) & ¿igual a $h_i$?\\
\midrule
0 & \texttt{00000000000000000000000000000000} & sí\\
1 & \texttt{692d1a9c017dc9fea6251bf243385b4b} & sí\\
2 & \texttt{a1815cf55fe8adaf93fbe63829aef30b} & no\\
3 & \texttt{d23ebb39f7eb3a5b140177367483078c} & no\\
4 & \texttt{587819e303f7bb5c105313742441de07} & no\\
\bottomrule
\end{tabular}
\end{center}

El verificador responde \texttt{SELLO INCORRECTO. Este registro NO es el publicado.}
En CBC el paso siguiente era ``recalcular todo y ya''. Aquí recalcular es igual de fácil
---no hay clave--- pero no sirve, porque el resultado tiene que ser igual a $S$, y $S$ no lo controlamos.
Las vías para cerrar el sello y el punto exacto en que cada una se detiene:

\begin{enumerate}[itemsep=3pt]
  \item \emph{Cambiar $S$.} Imposible: está publicado y replicado fuera del registro.
  \item \emph{Volver a la cadena original desde $h_2^*$.} Habría que encontrar $m_2^*$ con formato válido y
        $f(m_2^*,h_1)=h_2$, es decir una segunda preimagen de $f$; pero $m_2^*$ ya está fijado por lo que
        queremos que diga, y $f(m_2^*,h_1)=h_2^*\neq h_2$.
  \item \emph{Corregir con el último bloque.} Conocido $h_3^*$, buscar $m_4'$ con $E_{m_4'}(h_3^*)\xor h_3^* = S$, o sea
        \[
          E_{m_4'}(\texttt{d23ebb39\ldots83078c}) = S\xor h_3^* = \texttt{83f1099e4880ad609d9d541d3a485590}.
        \]
        \textbf{Aquí nos detenemos.} Es encontrar la clave AES que lleva un bloque conocido a otro bloque conocido:
        no hay despeje, solo búsqueda. Probamos 200\,000 claves de formato \texttt{TX04:L0177:\ldots} en 0,99\,s
        (unas $2\times10^5$ claves/s): ninguna. A ese ritmo, $2^{128}$ claves serían $\approx 5\times 10^{25}$ años.
        Peor: si $m_4'$ debe cumplir el formato, hay solo $\approx 1{,}8\times10^{10}$ candidatos, y el número esperado
        de soluciones es $1{,}8\times10^{10}/2^{128}\approx 5\times10^{-29}$: casi seguro \emph{no existe} ninguna.
  \item \emph{Retroceder desde el sello.} Para ir hacia atrás desde $S$ con $m_4$ habría que hallar $x=h_3$ tal que
        $E_{m_4}(x)\xor x = S$. Es la ecuación del literal (b)3 y tampoco admite despeje.
\end{enumerate}

\paso{Qué información faltaría para seguir}
Una forma de invertir $f$ en alguna de sus dos entradas: o bien una inversa de $x\mapsto E_m(x)\xor x$
(para retroceder desde $S$), o bien una forma de recuperar la clave de AES a partir de un par
(texto, cifrado) (para fabricar el último bloque). Ninguna de las dos existe; y nótese que no hay ningún
secreto que robar que lo resuelva: no hay clave.

\subsection*{(b) De qué depende}

\paso{1. Sin el término final}
\tesis{Sea $h_i = g(m_i,h_{i-1}) = E_{m_i}(h_{i-1})$. Conocidos $h_i\in\B$ y $m_i\in\mathcal{K}$, existe un único
$h_{i-1}\in\B$ que satisface la ecuación, y es calculable: $h_{i-1} = D_{m_i}(h_i)$.}

\paso{Hipótesis}
(H1) $m_i\in\mathcal{K}$ y $h_i\in\B$ son conocidos.
(H2) Para toda clave $m\in\mathcal K$, $E_m:\B\to\B$ es biyectiva, con inversa $D_m$ computable
(AES se descifra con la misma eficiencia con que se cifra).

\paso{Planteamiento}
Partimos de la ecuación $h_i = E_{m_i}(x)$ con incógnita $x\in\B$ y hay que probar dos cosas:
existencia de solución y unicidad, más que la solución se puede calcular.

\paso{Comparación}
\emph{Existencia.} Sea $x_0 := D_{m_i}(h_i)$, que está definido porque $D_{m_i}$ es una función en todo $\B$.
Entonces $E_{m_i}(x_0) = E_{m_i}(D_{m_i}(h_i)) = h_i$ por (H2), luego $x_0$ es solución.\\
\emph{Unicidad.} Sean $x, x'\in\B$ soluciones: $E_{m_i}(x) = h_i = E_{m_i}(x')$. Aplicando $D_{m_i}$ a ambos lados
y usando $D_{m_i}\circ E_{m_i} = \mathrm{id}$: $x = x'$. (Es la inyectividad de $E_{m_i}$.)\\
\emph{Despeje.} De $h_i = E_{m_i}(h_{i-1})$, aplicando $D_{m_i}$: $D_{m_i}(h_i) = D_{m_i}(E_{m_i}(h_{i-1})) = h_{i-1}$.\\
\emph{Calculabilidad.} $D_{m_i}(h_i)$ es un descifrado AES: una sola operación de coste constante.

\paso{Cierre lógico}
La solución existe (sobreyectividad), es única (inyectividad) y su fórmula es un descifrado AES, que se ejecuta
con los datos de (H1). Ninguno de los pasos requiere información adicional a $m_i$ y $h_i$.

\paso{Conclusión}
Sin el término final, cada paso de la cadena se deshace con $h_{i-1} = D_{m_i}(h_i)$. \hfill$\blacksquare$

Lo comprobamos en \texttt{experimentos.py}: para la variante sin término final, partiendo del sello
\texttt{44ee89ed\ldots eada27} y descifrando con $m_4, m_3, m_2, m_1$ se recuperan
$h_3, h_2, h_1$ y llegamos exactamente a $h_0 = 0$.

\paso{2. ¿Quién conoce $m_i$?}
Todo el mundo. En una cadena de bloques las entradas son públicas ---en \texttt{registro\_dm.json} están en
claro--- y \emph{tienen} que serlo: sin ellas nadie podría verificar recomputando. Así, en
$h_i=E_{m_i}(h_{i-1})$ la ``clave'' de la permutación es un dato público; la permutación entera, con su
inversa, está en manos de cualquiera. En CBC la reversibilidad estaba restringida a quien tuviera $k$;
aquí se le regala a todos.

Conclusión: \textbf{cualquiera} podría retroceder esa cadena, y \textbf{no necesita ningún secreto}.
Lo demostramos con una falsificación concreta: tomando la historia falsa (TX02 en EQ-A) y retrocediendo desde
el sello con $D_{m_4}, D_{m_3}, D_{m_2^*}, D_{m_1}$ obtuvimos
$h_0' = \texttt{4d311f3f3b922efe3b58a6d210b2bb83}$, y la historia falsa encadenada desde $h_0'$ da
exactamente el mismo sello. Si el punto de partida no está fijado, la falsificación es instantánea.
Si $h_0$ también está fijado, retroceder no basta por sí solo (hay que empalmar la parte que avanza desde $h_0$ con
la que retrocede desde $S$), pero el empalme es un encuentro a mitad de camino: variando $2^{64}$ bloques
en cada mitad se espera una coincidencia, $2^{64}$ operaciones en vez de $2^{128}$, y sigue sin hacer falta
ningún secreto. Es decir, sin el término final la seguridad de la cadena se derrumba.

\paso{3. Con el término final}
Para retroceder un paso habría que resolver, para la incógnita $x=h_{i-1}$,
\[
  E_{m_i}(x)\xor x = h_i .
\]
Intentamos el mismo despeje. Sumar $x$ en ambos lados da $E_{m_i}(x) = h_i\xor x$, y aplicar $D_{m_i}$ da
\[
  x = D_{m_i}(h_i\xor x),
\]
con $x$ todavía a los dos lados. Cada movimiento algebraico deja la incógnita en dos sitios: una vez
\emph{dentro} de la permutación no lineal $E_{m_i}$ y otra vez \emph{fuera}, sumada. No hay una operación que
saque a $x$ de ambos sitios a la vez. La función $x\mapsto E_m(x)\xor x$ no es una permutación (es la suma de una
permutación pseudoaleatoria y la identidad, y se comporta como una función aleatoria: tiene colisiones y valores sin
preimagen), así que ni siquiera hay garantía de que la solución exista o sea única, y no se conoce un método mejor que
probar valores de $x$: del orden de $2^{128}$ intentos.

\emph{Comparación con el Taller 0.} En $a\xor x = b$ la incógnita aparece \textbf{una sola vez}, y la operación que la
envuelve, sumar $a$, es una traslación del grupo $(\B,\xor)$ que se deshace sumando $a$ otra vez:
$x = a\xor b$ por (X1) y (X2). La ecuación de Davies--Meyer tiene lo que aquella no tenía: \textbf{la incógnita aparece dos
veces, una dentro de una permutación no lineal y otra fuera de ella}. El término $\xor h_{i-1}$, que parecía sobrar,
es exactamente lo que vuelve el paso no invertible aunque $E_{m_i}$ y $D_{m_i}$ sean públicas.

\subsection*{(c) Cerrar el caso}

\begin{center}
\renewcommand{\arraystretch}{1.3}
\begin{tabularx}{\linewidth}{>{\raggedright\arraybackslash}p{3.0cm}>{\raggedright\arraybackslash}X>{\raggedright\arraybackslash}X}
\toprule
 & \textbf{CBC} $c_i = E_k(m_i\xor c_{i-1})$ & \textbf{Davies--Meyer} $h_i = E_{m_i}(h_{i-1})\xor h_{i-1}$\\
\midrule
¿Recurrencia reversible? & Sí, para quien tiene $k$: $m_i = D_k(c_i)\xor c_{i-1}$ (biyección, ej.~2(b)). & No, para nadie: $E_{m_i}(x)\xor x=h_i$ no tiene despeje (ej.~3(b)3). Sin el término final sí lo sería, y para todos (ej.~3(b)1--2).\\
¿Qué hace falta para reescribir la historia? & La clave $k$: 4 cifrados AES, 22\,$\mu$s. Para borrar la primera o la última entrada, ni eso. & Invertir la compresión o hallar una clave AES para un par dado: $\approx2^{128}$ trabajo, y con formato válido casi seguro no hay solución. O bien sustituir el sello, que debe estar fuera de su alcance.\\
¿Qué protege? & Confidencialidad de las entradas frente a quien no tiene $k$; detecta corrupción accidental (vía formato). No da integridad. & Integridad de la historia completa frente al sello publicado. No da confidencialidad: las entradas están en claro.\\
¿Contra qué adversario? & Contra un tercero sin clave que quiera \emph{leer}. No contra el custodio (tiene $k$), ni contra un tercero que borre extremos o voltee bits. & Contra cualquiera, \emph{incluido el custodio}, siempre que el sello esté publicado fuera de su control y la historia se compare contra él.\\
\bottomrule
\end{tabularx}
\end{center}

\paso{¿Y si el sello estuviera dentro de \texttt{registro\_dm.json}?}
Entonces el mercado sí saldría impune. Procedimiento exacto:
\begin{enumerate}[itemsep=1pt]
  \item Editar la entrada: $m_2 \leftarrow m_2^* = \texttt{TX02:L0042:EQ-A\textvisiblespace}$.
  \item Recomputar hacia delante desde $h_0$ con la recurrencia, que no requiere clave:
        $h_1^*, h_2^*, h_3^*, h_4^* = \texttt{587819e303f7bb5c105313742441de07}$.
  \item Escribir ese $h_4^*$ como el nuevo sello dentro del archivo.
  \item El verificador recomputa, obtiene $h_4^*$, lo compara con el sello del propio archivo y responde ``sello correcto''.
\end{enumerate}
Lo hicimos (\texttt{dm\_con\_sello\_interno.json}, idéntico al \texttt{atk.json} recibido): un verificador que leyera el
sello del archivo lo acepta; contra el sello publicado, se rechaza. El cálculo cuesta lo mismo que una verificación
honesta. La función de compresión garantiza que \emph{a una historia le corresponde un único sello} en la práctica;
no garantiza \emph{cuál es el sello correcto}. Eso lo da el ancla.

Para que un registro sea realmente inmutable hace falta, además de la compresión no invertible:
\begin{itemize}[itemsep=1pt]
  \item un \textbf{ancla externa}: el sello publicado y replicado donde el custodio no pueda cambiarlo (otros nodos,
        una publicación con fecha, un servicio de sellado de tiempo, los propios usuarios guardando copia);
  \item \textbf{publicación oportuna}: el sello de la historia que incluye una transacción debe salir antes de que
        nadie tenga motivo para alterarla, y quien compra debe poder quedarse con él;
  \item fijar públicamente el punto de partida $h_0$ y la longitud de la historia, para que no se pueda empezar
        en otro sitio ni truncar.
\end{itemize}

%======================================================================
\section{La frase que hay que corregir}

La frase ``una cadena de bloques usa CBC'' es engañosa porque confunde \emph{la forma de la recurrencia} con
\emph{la propiedad que la hace útil}. Las dos fórmulas encadenan, sí, pero CBC encadena con una permutación
reversible para quien tiene la clave, y eso la hace un mecanismo de confidencialidad: el custodio puede descifrar,
editar y volver a cifrar sin dejar rastro. Una cadena de bloques necesita lo contrario: que nadie, ni siquiera
quien la lleva, pueda retroceder.

\textbf{Formulación correcta:} una cadena de bloques encadena cada bloque con el \emph{hash} del anterior, calculado con
una función de compresión resistente a preimágenes y colisiones (como Davies--Meyer sobre un cifrado de bloque), que
no se descifra sino que se verifica recomputando; y su inmutabilidad depende además de que el valor final esté anclado
en un lugar público que el custodio no controla.

%======================================================================
\section{Informe final}

\subsection{Las demostraciones}
Están en los ejercicios 2(b)1 (biyectividad de CBC) y 3(b)1 (reversibilidad de la variante sin término final),
cada una con sus pasos rotulados: Hipótesis, Planteamiento, Comparación, Cierre lógico y Conclusión, y con la tesis
enunciada en recuadro. El análisis de 3(b)3 explica por qué no hay demostración análoga para Davies--Meyer.

\subsection{Código}
\subsubsection*{\texttt{registro\_tool.py} (con las tres funciones completadas)}
\lstinputlisting[style=py]{../codigo/registro_tool.py}

\subsubsection*{\texttt{experimentos.py}}
\lstinputlisting[style=py]{../codigo/experimentos.py}

\subsection{Cadena falsificada y salida del verificador}
Archivo \texttt{falsificado.json} (ejercicio 2; el campo \texttt{cadena}, que en el archivo es una sola
línea de 128 caracteres hexadecimales, se parte aquí en sus cuatro bloques):
\begin{lstlisting}[style=salida]
{
  "modo": "CBC",
  "formula": "c_i = E_k(m_i XOR c_{i-1})",
  "clave": "4ac43a52cc764438cae7b3af9fced5a9",
  "iv": "3c1cc2d73a499ca72f36e7303d76a130",
  "cadena": "acc42ecaebdb06d4e3b2852697a91947
             8f3ebfc015f257f20935a4a9bb070b22
             3314b26facf1eaae43ea1b86beb9ec3b
             0b1fc3366cb2c953363114fc04286dc5"
}
\end{lstlisting}
Salida literal de \texttt{python3 verificador.py cbc falsificado.json}:
\lstinputlisting[style=salida]{../codigo/salida_verificador_cbc_falsificado.txt}

Para contraste, la salida de \texttt{python3 verificador.py dm dm\_falsificado.json} (la misma alteración sobre
Davies--Meyer):
\lstinputlisting[style=salida]{../codigo/salida_verificador_dm_falsificado.txt}

\subsection{Tabla comparativa}
Ver la tabla del ejercicio 3(c).

\subsection{Dictamen para el usuario}

\emph{Estimado usuario:}

Usted preguntó si el mercado pudo haber borrado su compra. La respuesta corta es: \textbf{sí, pudo}, y
con el sistema que usa hoy nadie lo habría notado.

El mercado guarda el registro en una especie de caja fuerte encadenada, en la que cada transferencia queda
atada a la anterior. Eso suena muy seguro, pero el candado de esa caja tiene una sola llave, y la llave la tiene el
propio mercado. Con ella puede abrir el registro, cambiar o quitar lo que quiera, volver a cerrarlo y dejarlo
exactamente igual de ``encadenado'' que antes. Nosotros lo hicimos: cambiamos una transferencia en una fracción de
segundo y la herramienta oficial del mercado dijo que todo estaba en orden. Incluso sin la llave se puede quitar la
primera o la última transferencia sin que la herramienta proteste.

Por eso, cuando el mercado dice ``la cadena verifica, así que no pudo pasar'', está diciendo algo que no es cierto. Que
la herramienta diga que todo está bien no demuestra que nadie haya tocado el registro. Esto no prueba que el mercado
le haya borrado su compra; prueba que su defensa no vale.

Para que la respuesta fuera ``no, imposible'', el mercado tendría que cambiar dos cosas:
\begin{enumerate}[itemsep=1pt]
  \item Dejar de usar un candado con llave y usar un ``sello'': una huella del registro completo que cualquiera puede
        recalcular, que cambia por completo si se toca una sola letra, y que nadie ---ni siquiera el mercado--- sabe
        fabricar para un registro distinto.
  \item Publicar esa huella en un sitio que el mercado no controle (por ejemplo, entregársela a cada comprador en su
        recibo, o publicarla en varios lugares independientes). Si el mercado guarda la huella en su propio cajón, puede
        cambiar el registro y la huella a la vez.
\end{enumerate}
Con esas dos condiciones, usted guardaría la huella del día de su compra y, si alguien borrara su transferencia, la
huella dejaría de coincidir y cualquiera podría comprobarlo.

\subsection{Media página por integrante: por qué ``una cadena de bloques usa CBC'' es engañoso}

\emph{Nota: cada integrante debe redactar o revisar su propia sección. Lo que sigue son cuatro enfoques distintos
como punto de partida.}

\subsubsection*{Integrante 1: \underline{\hspace{4cm}} --- La forma no es la propiedad}
Las fórmulas $c_i = E_k(m_i\xor c_{i-1})$ y $h_i = H(B_i, h_{i-1})$ se parecen porque las dos son recurrencias: cada
término depende del anterior. Pero ``depender del anterior'' es una propiedad de la forma, y la seguridad no está en
la forma sino en la función que se aplica en cada paso. En CBC esa función es una permutación que, con la clave, se
deshace; demostramos que la aplicación historia$\to$cadena es una biyección con inversa explícita. Un objeto
biyectivo y calculable en ambos sentidos no puede servir de prueba de nada: cada historia tiene su cadena y cada
cadena su historia, y pasar de una a otra cuesta microsegundos. Decir que una cadena de bloques ``usa CBC'' es tomar
el parecido de las fórmulas por un parecido de propiedades. Lo correcto: una cadena de bloques encadena con una
función de compresión no invertible, y se verifica recomputando, no descifrando.

\subsubsection*{Integrante 2: \underline{\hspace{4cm}} --- ¿Quién tiene la llave?}
En CBC la única barrera para reescribir la historia es la clave, y en un registro la clave la tiene precisamente quien
lo custodia: el adversario del que queremos protegernos. CBC protege contra terceros que quieren \emph{leer}, no contra
el custodio que quiere \emph{editar}. Además mostramos que, aun sin la clave, borrar la primera o la última entrada
deja una cadena que el verificador acepta. Una cadena de bloques invierte el planteamiento: no hay llave, las entradas
son públicas y la protección no consiste en ocultar sino en que nadie pueda fabricar una historia distinta con el mismo
sello. Por eso la formulación correcta no menciona cifrado: una cadena de bloques es una cadena de hashes, cuya
seguridad descansa en la dificultad de invertir la compresión, y no en un secreto.

\subsubsection*{Integrante 3: \underline{\hspace{4cm}} --- Un solo término}
El ejercicio 3 muestra que la diferencia puede ser un solo término. Davies--Meyer sin el $\xor h_{i-1}$ final es
$h_i = E_{m_i}(h_{i-1})$, una permutación de $h_{i-1}$ cuya ``clave'' es el mensaje, que es público; cualquiera la
retrocede con $D_{m_i}$. Esa variante sí ``se parece a CBC'' en lo esencial ---es reversible--- y es insegura por eso.
Con el término final, la ecuación $E_{m_i}(x)\xor x = h_i$ tiene la incógnita dentro y fuera de una permutación y no se
despeja. La moraleja es que lo que hace a una cadena de bloques no es usar un cifrado de bloque (Davies--Meyer usa AES),
sino usarlo de forma que el resultado no sea invertible. Decir ``usa CBC'' esconde justo el detalle que importa.

\subsubsection*{Integrante 4: \underline{\hspace{4cm}} --- El ancla}
Incluso con una función no invertible, el ejercicio 3(c) muestra que si el sello vive dentro del archivo, el custodio
cambia una entrada, recalcula y reemplaza el sello, sin ninguna clave. La función de compresión solo garantiza que a
cada historia le corresponde un sello; no dice cuál es el correcto. Eso lo decide un valor publicado que el custodio no
controla. En una cadena de bloques ese papel lo cumplen la replicación y el consenso entre muchos participantes. Por
eso la frase ``una cadena de bloques usa CBC'' es doblemente engañosa: atribuye la seguridad a un modo de cifrado,
cuando proviene de dos cosas que CBC no tiene: una compresión que no se deja invertir y un ancla pública fuera del
alcance de quien lleva el registro.

\end{document}
